% Fichier complété par tools/complete_missing_corrections.py.
% Source : PPM/PPM-03/85_FauteuilBateau/85_FauteuilBateau.tex
% Statut : rédigé (6 question(s)).
\begin{corrigebox}[Corrigé rédigé]
\CorrigeQuestion{1}
Les surfaces de référence sont choisies parce qu'elles
                réalisent la mise en position fonctionnelle de la pièce dans l'assemblage.
                La référence primaire supprime trois degrés de liberté, la secondaire deux
                et la tertiaire le dernier; elles doivent être étendues, stables,
                accessibles au contrôle et directement liées aux surfaces d'appui réelles.
\CorrigeQuestion{2}
La spécification se décode en identifiant successivement
                l'élément tolérancé, le symbole géométrique, la valeur et la forme de la
                zone de tolérance, les références ordonnées et les modificateurs. Le
                gabarit est le simulateur géométrique associé : deux plans pour une forme
                ou une orientation de surface, un cylindre pour un axe, ou une zone
                cylindrique/sphérique pour une localisation. La pièce est conforme si
                l'élément extrait reste entièrement dans cette zone après mise en référence.
\CorrigeQuestion{3}
La spécification se décode en identifiant successivement
                l'élément tolérancé, le symbole géométrique, la valeur et la forme de la
                zone de tolérance, les références ordonnées et les modificateurs. Le
                gabarit est le simulateur géométrique associé : deux plans pour une forme
                ou une orientation de surface, un cylindre pour un axe, ou une zone
                cylindrique/sphérique pour une localisation. La pièce est conforme si
                l'élément extrait reste entièrement dans cette zone après mise en référence.
\CorrigeQuestion{4}
Le modificateur au maximum de matière autorise une
                tolérance bonus égale à l'écart entre la taille réelle et la taille au
                maximum de matière. Sur l'élément tolérancé, la zone de position s'agrandit
                lorsque l'on s'éloigne du maximum de matière; appliqué à la référence, il
                crée une mobilité de datum correspondant au jeu disponible dans le
                simulateur fonctionnel.
\CorrigeQuestion{5}
On choisit un brut proche de la forme (forge, moulage ou
                débit de barre), puis : création des références, ébauche de toutes les
                surfaces, stabilisation/traitement éventuel, finition des alésages et
                portées en reprenant les références fonctionnelles, perçages et taraudages,
                rectification si nécessaire, ébavurage, traitement de surface et contrôle
                GPS final. Chaque phase utilise une mise en position 3--2--1 cohérente avec
                le système de références du plan.
\CorrigeQuestion{6}
On choisit un brut proche de la forme (forge, moulage ou
                débit de barre), puis : création des références, ébauche de toutes les
                surfaces, stabilisation/traitement éventuel, finition des alésages et
                portées en reprenant les références fonctionnelles, perçages et taraudages,
                rectification si nécessaire, ébavurage, traitement de surface et contrôle
                GPS final. Chaque phase utilise une mise en position 3--2--1 cohérente avec
                le système de références du plan.
\end{corrigebox}
